Объект разработки: программный комплекс для интерполяции многомерных отображений и решений дифференциальных уравнений на основе адаптивных разреженных сеток.

Цель работы: реализация и исследование параллельных алгоритмов интерполяции на основе разреженных сеток с квадратичным базисом, а также создание прикладной системы для удалённого вычисления и визуализации результатов.

Методы и методология: методы численного анализа и интерполяции; иерархические базисные функции; адаптивное уточнение по коэффициентам избытка; использование контрольных точек; численное интегрирование систем обыкновенных дифференциальных уравнений методом Рунге--Кутты--Фельберга; методы параллельного программирования; клиент-серверная архитектура с использованием gRPC и protobuf.

Результаты работы и их новизна: разработано вычислительное ядро на языке Go, реализующее построение адаптивной разреженной сетки для скалярных и векторных отображений в последовательном и параллельном режимах. Реализована поддержка квадратичного базиса, позволяющего получать более гладкую аппроксимацию по сравнению с линейным вариантом. Создан серверный слой удалённого доступа, интерпретирующий пользовательские формулы и поддерживающий аппроксимацию операторов, возникающих при решении дифференциальных уравнений. Разработано клиентское приложение визуализации на Flutter. Научная новизна состоит в объединении адаптивной разреженной интерполяции, параллельного построения и прикладной инфраструктуры визуального анализа в рамках единого программного комплекса.

Область применения: визуализация и исследование двумерных динамических систем; численный эксперимент с операторами перехода; демонстрация методов адаптивной интерполяции в учебных курсах; использование как основы для дальнейших исследований в области параллельных численных методов.